\documentclass[../../main.tex]{subfiles}% 注意这里的文件路径不能用 ./main.tex ,否则用latexmk编译子文件会报错
\graphicspath{{\subfix{./image/}}{\subfix{../../image/}}} % 指定图片引用路径,后续可以直接使用图片文件名
% 如果图片存放在上级目录,则使用 ../../image/ 作为备用路径

% 例如:
% \begin{figure}[H]
% \centering
% \includegraphics[width=0.3\textwidth]{图.png}
% \caption{}
% \label{figure:图}
% \end{figure}
% 注意:上述\label{}一定要放在\caption{}之后,否则引用图片序号会只会显示??.

\begin{document}

% SOURCE p=39; printed=24; object=ch01-sec03-title; region=lower
\section{测度}\label{section:ch01-measures}

% SOURCE p=39; printed=24; object=ch01-measure-definition; region=bottom

\begin{definition}[测度与可数可加性]\label{definition:ch01-measure-definition}
设 $X$ 为配备了 $\sigma$-代数 $\mathscr{M}$ 的集合。$\mathscr{M}$ 上（或 $(X,\mathscr{M})$ 上，或在 $\mathscr{M}$ 已明确时简称 $X$ 上）的\textbf{测度}，是满足下列条件的函数 $\mu:\mathscr{M}\to[0,\infty]$：
\begin{enumerate}[label=\roman*.,ref=(\roman*)]
\item\label{item:ch01-measure-empty} $\mu(\varnothing)=0$；
\item\label{item:ch01-measure-countably-additive} 若 $\{E_j\}_1^\infty$ 是 $\mathscr{M}$ 中两两不交的集合序列，则 $\mu(\bigcup_1^\infty E_j)=\sum_1^\infty\mu(E_j)$。
\end{enumerate}
\cref{item:ch01-measure-countably-additive}称为\textbf{可数可加性}。
\end{definition}
% SOURCE p=40; printed=25; object=ch01-finitely-additive-definition; region=top

\begin{definition}[有限可加性]\label{definition:ch01-finite-additivity}
下面的性质称为\textbf{有限可加性}：
\begin{enumerate}[label=\roman*$'$.,ref=(\roman*$'$),start=2]
\item\label{item:ch01-measure-finitely-additive} 若 $E_1,\ldots,E_n$ 是 $\mathscr{M}$ 中两两不交的集合，则 $\mu(\bigcup_1^nE_j)=\sum_1^n\mu(E_j)$。
\end{enumerate}
\end{definition}
\begin{proposition}\label{proposition:ch01-countable-implies-finite-additivity}
可数可加性蕴含有限可加性。
\end{proposition}
\begin{proof}
因为可以对 $j>n$ 取 $E_j=\varnothing$。
\end{proof}
\begin{definition}[有限可加测度]\label{definition:ch01-finitely-additive-definition}
满足\cref{item:ch01-measure-empty,item:ch01-measure-finitely-additive}、但未必满足\cref{item:ch01-measure-countably-additive}的函数 $\mu$ 称为\textbf{有限可加测度}。
\end{definition}


% SOURCE p=40; printed=25; object=ch01-measure-space-definition; region=upper

\begin{definition}[可测空间与测度空间]\label{definition:ch01-measure-space-definition}
若 $X$ 为集合且 $\mathscr{M}\subseteq\mathscr{P}(X)$ 为 $\sigma$-代数，则 $(X,\mathscr{M})$ 称为\textbf{可测空间}，$\mathscr{M}$ 中的集合称为\textbf{可测集}。若 $\mu$ 为 $(X,\mathscr{M})$ 上的测度，则 $(X,\mathscr{M},\mu)$ 称为\textbf{测度空间}。
\end{definition}


% SOURCE p=40; printed=25; object=ch01-finiteness-definitions; region=upper-middle
设 $(X,\mathscr{M},\mu)$ 为测度空间。下面给出一些描述 $\mu$“大小”的标准术语。
\begin{definition}[有限、$\sigma$-有限与半有限测度]\label{definition:ch01-finiteness-definitions}
若 $\mu(X)<\infty$，则称 $\mu$ \textbf{有限}。（这蕴含对所有 $E\in\mathscr{M}$，$\mu(E)<\infty$，因为 $\mu(X)=\mu(E)+\mu(E^{\mathrm{c}})$。）若 $X=\bigcup_1^\infty E_j$，其中对所有 $j$ 都有 $E_j\in\mathscr{M}$ 且 $\mu(E_j)<\infty$，则称 $\mu$ 为 \textbf{$\sigma$-有限}的。更一般地，若 $E=\bigcup_1^\infty E_j$，其中对所有 $j$ 都有 $E_j\in\mathscr{M}$ 且 $\mu(E_j)<\infty$，则称集合 $E$ 关于 $\mu$ 是 \textbf{$\sigma$-有限}的。（说 $E$ 具有 $\sigma$-有限测度也正确，但更为冗长。）若对每个满足 $\mu(E)=\infty$ 的 $E\in\mathscr{M}$，存在 $F\in\mathscr{M}$ 使 $F\subseteq E$ 且 $0<\mu(F)<\infty$，则称 $\mu$ \textbf{半有限}。
\end{definition}


% SOURCE p=40; printed=25; object=ch01-finiteness-discussion; region=middle

\begin{remark}
每个 $\sigma$-有限测度都是半有限的（\cref{exercise:ch01-sigma-finite-semifinite}），但反之不成立。实际中出现的测度大多是 $\sigma$-有限的，这是幸运的，因为非 $\sigma$-有限测度往往呈现病态行为。习题中会不时探讨非 $\sigma$-有限测度的性质。
\end{remark}


% SOURCE p=40; printed=25; object=ch01-measure-examples-intro; region=middle
我们来看几个测度的例子。这些例子都相当简单，不过第一个具有实际重要性。构造更有趣的例子，将是随后两节的任务。

% SOURCE p=40; printed=25; object=ch01-measure-examples; region=lower

\begin{examples}\label{examples:ch01-measure-examples}
\begin{itemize}
\item 设 $X$ 为任意非空集合，$\mathscr{M}=\mathscr{P}(X)$，$f$ 为从 $X$ 到 $[0,\infty]$ 的任意函数。则 $f$ 通过公式 $\mu(E)=\sum_{x\in E}f(x)$ 确定 $\mathscr{M}$ 上的一个测度 $\mu$。（这种可能不可数的和的定义见\cref{section:ch00-extended-real-numbers}。）读者可验证，$\mu$ 半有限当且仅当对每个 $x\in X$ 都有 $f(x)<\infty$；$\mu$ 为 $\sigma$-有限测度，当且仅当它半有限且 $\{x:f(x)>0\}$ 可数。
\begin{definition*}[计数测度与 Dirac 测度]
两个特例尤为重要：若对所有 $x$ 都有 $f(x)=1$，则 $\mu$ 称为\textbf{计数测度}；若对于某个 $x_0\in X$，定义 $f(x_0)=1$，并在 $x\ne x_0$ 时定义 $f(x)=0$，则 $\mu$ 称为 $x_0$ 处的\textbf{点质量}或 \textbf{Dirac 测度}。（这些名称也用于这些测度在 $X$ 上更小的 $\sigma$-代数上的限制。）
\end{definition*}

\item 设 $X$ 为不可数集合，$\mathscr{M}$ 为可数集或余可数集的 $\sigma$-代数。在 $\mathscr{M}$ 上定义函数 $\mu$：当 $E$ 可数时 $\mu(E)=0$，当 $E$ 余可数时 $\mu(E)=1$。容易看出 $\mu$ 为测度。
\item 设 $X$ 为无限集，$\mathscr{M}=\mathscr{P}(X)$。当 $E$ 有限时定义 $\mu(E)=0$，当 $E$ 无限时定义 $\mu(E)=\infty$。则 $\mu$ 为有限可加测度，但不是测度。
\end{itemize}
\end{examples}


% SOURCE p=40; printed=25; object=ch01-measure-properties-intro; region=lower
测度的基本性质归纳于下述定理。

% SOURCE p=40; printed=25; object=ch01-measure-properties-start; number=1.8; region=bottom
\begin{theorem}\label{theorem:ch01-measure-properties}
设 $(X,\mathscr{M},\mu)$ 为测度空间。
\begin{enumerate}[label=\alph*.,ref=(\alph*)]
\item\label{item:ch01-measure-monotonicity} （\textbf{单调性}）若 $E,F\in\mathscr{M}$ 且 $E\subseteq F$，则 $\mu(E)\leqslant\mu(F)$。
\item\label{item:ch01-measure-subadditivity} （\textbf{次可加性}）若 $\{E_j\}_1^\infty\subseteq\mathscr{M}$，则 $\mu(\bigcup_1^\infty E_j)\leqslant\sum_1^\infty\mu(E_j)$。
% SOURCE p=41; printed=26; object=ch01-measure-properties-end; region=top
\item\label{item:ch01-measure-continuity-below} （\textbf{下连续性}）若 $\{E_j\}_1^\infty\subseteq\mathscr{M}$ 且 $E_1\subseteq E_2\subseteq\cdots$，则 $\mu(\bigcup_1^\infty E_j)=\lim_{j\to\infty}\mu(E_j)$。
\item\label{item:ch01-measure-continuity-above} （\textbf{上连续性}）若 $\{E_j\}_1^\infty\subseteq\mathscr{M}$，$E_1\supseteq E_2\supseteq\cdots$，且 $\mu(E_1)<\infty$，则 $\mu(\bigcap_1^\infty E_j)=\lim_{j\to\infty}\mu(E_j)$。
\end{enumerate}
\end{theorem}

% SOURCE p=41; printed=26; object=ch01-proof-measure-monotonicity; region=upper
\begin{proof}
\subcref{theorem:ch01-measure-properties}{item:ch01-measure-monotonicity}：若 $E\subseteq F$，则 $\mu(F)=\mu(E)+\mu(F\setminus E)\geqslant\mu(E)$。

% SOURCE p=41; printed=26; object=ch01-proof-measure-subadditivity; region=upper
\subcref{theorem:ch01-measure-properties}{item:ch01-measure-subadditivity}：令 $F_1=E_1$，并对 $k>1$ 令 $F_k=E_k\setminus(\bigcup_1^{k-1}E_j)$。则各 $F_k$ 两两不交，且对所有 $n$，$\bigcup_1^nF_j=\bigcup_1^nE_j$。因此，由\subcref{theorem:ch01-measure-properties}{item:ch01-measure-monotonicity}，
\begin{align*}
\mu\Bigl(\bigcup_1^\infty E_j\Bigr)
=\mu\Bigl(\bigcup_1^\infty F_j\Bigr)
=\sum_1^\infty\mu(F_j)\leqslant\sum_1^\infty\mu(E_j).
\end{align*}

% SOURCE p=41; printed=26; object=ch01-proof-measure-continuity-below; region=middle
\subcref{theorem:ch01-measure-properties}{item:ch01-measure-continuity-below}：令 $E_0=\varnothing$，则
\begin{align*}
\mu\Bigl(\bigcup_1^\infty E_j\Bigr)
&=\sum_1^\infty\mu(E_j\setminus E_{j-1})\\
&=\lim_{n\to\infty}\sum_1^n\mu(E_j\setminus E_{j-1})=\lim_{n\to\infty}\mu(E_n).
\end{align*}

% SOURCE p=41; printed=26; object=ch01-proof-measure-continuity-above; region=middle
\subcref{theorem:ch01-measure-properties}{item:ch01-measure-continuity-above}：令 $F_j=E_1\setminus E_j$；则 $F_1\subseteq F_2\subseteq\cdots$，$\mu(E_1)=\mu(F_j)+\mu(E_j)$，且 $\bigcup_1^\infty F_j=E_1\setminus(\bigcap_1^\infty E_j)$。于是，由\subcref{theorem:ch01-measure-properties}{item:ch01-measure-continuity-below}，
\begin{align*}
\mu(E_1)&=\mu\Bigl(\bigcap_1^\infty E_j\Bigr)+\lim_{j\to\infty}\mu(F_j)\\
&=\mu\Bigl(\bigcap_1^\infty E_j\Bigr)+\lim_{j\to\infty}[\mu(E_1)-\mu(E_j)].
\end{align*}
由于 $\mu(E_1)<\infty$，可从两边减去它，即得所需结论。
\end{proof}

% SOURCE p=41; printed=26; object=ch01-continuity-finiteness-caveat; region=middle-lower

\begin{remark}
注意，\subcref{theorem:ch01-measure-properties}{item:ch01-measure-continuity-above}中的条件 $\mu(E_1)<\infty$ 可替换为对某个 $n>1$ 有 $\mu(E_n)<\infty$，因为可以从序列中去掉前 $n-1$ 个 $E_j$ 而不改变其交。不过，某种有限性假设是必要的，因为可能对所有 $j$ 都有 $\mu(E_j)=\infty$，却有 $\mu(\bigcap_1^\infty E_j)<\infty$。（例如，取 $\mu$ 为 $(\mathbb{N},\mathscr{P}(\mathbb{N}))$ 上的计数测度，令 $E_j=\{n:n\geqslant j\}$；则 $\bigcap_1^\infty E_j=\varnothing$。）
\end{remark}


% SOURCE p=41; printed=26; object=ch01-null-almost-everywhere; region=lower

\begin{definition}[零测集]\label{definition:ch01-null-almost-everywhere}
若 $(X,\mathscr{M},\mu)$ 为测度空间，则满足 $\mu(E)=0$ 的集合 $E\in\mathscr{M}$ 称为\textbf{零测集}。
\end{definition}

\begin{remark}
由次可加性，零测集的任意可数并仍为零测集；我们将经常使用这一事实。
\end{remark}

\begin{definition}[几乎处处]\label{definition:ch01-almost-everywhere}
若关于点 $x\in X$ 的一个断言除某个零测集中的 $x$ 外都成立，则称它\textbf{几乎处处}成立（缩写为 \textbf{a.e.}），或称它对\textbf{几乎每个} $x$ 成立。（需要更精确时，我们会说 \textbf{$\mu$-零测集}或 \textbf{$\mu$-几乎处处}。）
\end{definition}


% SOURCE p=41; printed=26; object=ch01-complete-measure-definition; region=lower
若 $\mu(E)=0$ 且 $F\subseteq E$，则只要 $F\in\mathscr{M}$，由单调性就有 $\mu(F)=0$；但一般未必有 $F\in\mathscr{M}$。
\begin{definition}[完备测度]\label{definition:ch01-complete-measure-definition}
若一个测度的定义域包含所有零测集的所有子集，则称该测度\textbf{完备}。
\end{definition}
完备性有时可以避免一些烦人的技术细节，而且总可以通过扩大 $\mu$ 的定义域来达到，如下所述。

% SOURCE p=41; printed=26; object=ch01-measure-completion; number=1.9; region=bottom
\begin{theorem}\label{theorem:ch01-measure-completion}
设 $(X,\mathscr{M},\mu)$ 为测度空间。令 $\mathscr{N}=\{N\in\mathscr{M}:\mu(N)=0\}$，并令
\begin{align*}
\overline{\mathscr{M}}=\{E\cup F:E\in\mathscr{M},\ \text{且存在 }N\in\mathscr{N}\text{ 使 }F\subseteq N\}.
\end{align*}
则 $\overline{\mathscr{M}}$ 是 $\sigma$-代数，且 $\mu$ 有唯一的延拓 $\overline{\mu}$，使其成为 $\overline{\mathscr{M}}$ 上的完备测度。
\end{theorem}

% SOURCE p=42; printed=27; object=ch01-proof-measure-completion-algebra; region=top
\begin{proof}
由于 $\mathscr{M}$ 和 $\mathscr{N}$ 对可数并封闭，$\overline{\mathscr{M}}$ 也如此。若 $E\cup F\in\overline{\mathscr{M}}$，其中 $E\in\mathscr{M}$ 且 $F\subseteq N\in\mathscr{N}$，则可设 $E\cap N=\varnothing$（否则，用 $F\setminus E$ 和 $N\setminus E$ 分别替换 $F$ 和 $N$）。于是 $E\cup F=(E\cup N)\cap(N^{\mathrm{c}}\cup F)$，所以 $(E\cup F)^{\mathrm{c}}=(E\cup N)^{\mathrm{c}}\cup(N\setminus F)$。但 $(E\cup N)^{\mathrm{c}}\in\mathscr{M}$ 且 $N\setminus F\subseteq N$，因而 $(E\cup F)^{\mathrm{c}}\in\overline{\mathscr{M}}$。故 $\overline{\mathscr{M}}$ 为 $\sigma$-代数。

% SOURCE p=42; printed=27; object=ch01-proof-measure-completion-extension; region=upper
若 $E\cup F\in\overline{\mathscr{M}}$ 如上所述，定义 $\overline{\mu}(E\cup F)=\mu(E)$。这一定义是良定义的，因为若 $E_1\cup F_1=E_2\cup F_2$，其中 $F_j\subseteq N_j\in\mathscr{N}$，则 $E_1\subseteq E_2\cup N_2$，所以 $\mu(E_1)\leqslant\mu(E_2)+\mu(N_2)=\mu(E_2)$；同理，$\mu(E_2)\leqslant\mu(E_1)$。容易验证，$\overline{\mu}$ 是 $\overline{\mathscr{M}}$ 上的完备测度，并且是 $\overline{\mathscr{M}}$ 上唯一延拓 $\mu$ 的测度；细节留给读者（\cref{exercise:ch01-completion-proof}）。
\end{proof}

% SOURCE p=42; printed=27; object=ch01-completion-terminology; region=upper-middle

\begin{definition}[测度及 $\sigma$-代数的完备化]\label{definition:ch01-completion-terminology}
\cref{theorem:ch01-measure-completion}中的测度 $\overline{\mu}$ 称为 $\mu$ 的\textbf{完备化}，$\overline{\mathscr{M}}$ 称为 $\mathscr{M}$ 关于 $\mu$ 的\textbf{完备化}。
\end{definition}


% SOURCE p=42; printed=27; object=ch01-sec03-exercises; region=middle
\subsection*{习题}

% SOURCE p=42; printed=27; object=ch01-exercise-completion-proof; number=6; region=middle
\begin{exercise}\label{exercise:ch01-completion-proof}
补全\cref{theorem:ch01-measure-completion}的证明。
\end{exercise}

% SOURCE p=42; printed=27; object=ch01-exercise-linear-combination; number=7; region=middle
\begin{exercise}\label{exercise:ch01-linear-combination}
若 $\mu_1,\ldots,\mu_n$ 为 $(X,\mathscr{M})$ 上的测度，且 $a_1,\ldots,a_n\in[0,\infty)$，则 $\sum_1^na_j\mu_j$ 是 $(X,\mathscr{M})$ 上的测度。
\end{exercise}

% SOURCE p=42; printed=27; object=ch01-exercise-set-limits; number=8; region=middle
\begin{exercise}\label{exercise:ch01-set-limits}
若 $(X,\mathscr{M},\mu)$ 为测度空间，且 $\{E_j\}_1^\infty\subseteq\mathscr{M}$，则 $\mu(\liminf E_j)\leqslant\liminf\mu(E_j)$。此外，若 $\mu(\bigcup_1^\infty E_j)<\infty$，则 $\mu(\limsup E_j)\geqslant\limsup\mu(E_j)$。
\end{exercise}

% SOURCE p=42; printed=27; object=ch01-exercise-inclusion-exclusion; number=9; region=middle
\begin{exercise}\label{exercise:ch01-inclusion-exclusion}
若 $(X,\mathscr{M},\mu)$ 为测度空间，且 $E,F\in\mathscr{M}$，则 $\mu(E)+\mu(F)=\mu(E\cup F)+\mu(E\cap F)$。
\end{exercise}

% SOURCE p=42; printed=27; object=ch01-exercise-restriction; number=10; region=middle-lower
\begin{exercise}\label{exercise:ch01-restriction}
给定测度空间 $(X,\mathscr{M},\mu)$ 及 $E\in\mathscr{M}$，对 $A\in\mathscr{M}$ 定义 $\mu_E(A)=\mu(A\cap E)$。则 $\mu_E$ 是测度。
\end{exercise}

% SOURCE p=42; printed=27; object=ch01-exercise-continuity-additivity; number=11; region=lower
\begin{exercise}\label{exercise:ch01-continuity-additivity}
有限可加测度 $\mu$ 是测度，当且仅当它具有\subcref{theorem:ch01-measure-properties}{item:ch01-measure-continuity-below}中的下连续性。若 $\mu(X)<\infty$，则 $\mu$ 是测度，当且仅当它具有\subcref{theorem:ch01-measure-properties}{item:ch01-measure-continuity-above}中的上连续性。
\end{exercise}

% SOURCE p=42; printed=27; object=ch01-exercise-measure-algebra-metric; number=12; region=lower
\begin{exercise}\label{exercise:ch01-measure-algebra-metric}
设 $(X,\mathscr{M},\mu)$ 为有限测度空间。
\begin{enumerate}[label=\alph*.,ref=(\alph*)]
\item\label{item:ch01-symmetric-difference-zero} 若 $E,F\in\mathscr{M}$ 且 $\mu(E\mathbin{\triangle}F)=0$，则 $\mu(E)=\mu(F)$。
\item\label{item:ch01-symmetric-difference-equivalence} 规定 $E\sim F$ 当且仅当 $\mu(E\mathbin{\triangle}F)=0$；则 $\sim$ 是 $\mathscr{M}$ 上的等价关系。
\item\label{item:ch01-symmetric-difference-metric} 对 $E,F\in\mathscr{M}$，定义 $\rho(E,F)=\mu(E\mathbin{\triangle}F)$。则 $\rho(E,G)\leqslant\rho(E,F)+\rho(F,G)$，从而 $\rho$ 在等价类空间 $\mathscr{M}/{\sim}$ 上定义了一个度量。
\end{enumerate}
\end{exercise}

% SOURCE p=42; printed=27; object=ch01-exercise-sigma-finite-semifinite; number=13; region=lower
\begin{exercise}\label{exercise:ch01-sigma-finite-semifinite}
每个 $\sigma$-有限测度都是半有限的。
\end{exercise}

% SOURCE p=42; printed=27; object=ch01-exercise-semifinite-large-finite; number=14; region=bottom
\begin{exercise}\label{exercise:ch01-semifinite-large-finite}
若 $\mu$ 为半有限测度且 $\mu(E)=\infty$，则对任意 $C>0$，存在 $F\subseteq E$ 使 $C<\mu(F)<\infty$。
\end{exercise}

% SOURCE p=42; printed=27; object=ch01-exercise-semifinite-part-start; number=15; region=bottom
\begin{exercise}\label{exercise:ch01-semifinite-part}
给定 $(X,\mathscr{M})$ 上的测度 $\mu$，在 $\mathscr{M}$ 上定义 $\mu_0(E)=\sup\{\mu(F):F\subseteq E\text{ 且 }\mu(F)<\infty\}$。
\begin{enumerate}[label=\alph*.,ref=(\alph*)]
\item\label{item:ch01-semifinite-part-measure} $\mu_0$ 是半有限测度，称为 $\mu$ 的\textbf{半有限部分}。
\item\label{item:ch01-semifinite-part-equality} 若 $\mu$ 半有限，则 $\mu=\mu_0$。（使用\cref{exercise:ch01-semifinite-large-finite}。）
% SOURCE p=43; printed=28; object=ch01-exercise-semifinite-part-end; region=top
\item\label{item:ch01-semifinite-part-remainder} 存在 $\mathscr{M}$ 上只取 $0$ 和 $\infty$ 两个值的测度 $\nu$（一般不唯一），使得 $\mu=\mu_0+\nu$。
\end{enumerate}
\end{exercise}

% SOURCE p=43; printed=28; object=ch01-exercise-saturation; number=16; region=upper
\begin{exercise}\label{exercise:ch01-saturation}

\begin{definition*}[局部可测集与饱和测度]
设 $(X,\mathscr{M},\mu)$ 为测度空间。集合 $E\subseteq X$ 称为\textbf{局部可测}的，是指对所有满足 $\mu(A)<\infty$ 的 $A\in\mathscr{M}$，都有 $E\cap A\in\mathscr{M}$。令 $\widetilde{\mathscr{M}}$ 为所有局部可测集构成的族。显然，$\mathscr{M}\subseteq\widetilde{\mathscr{M}}$；若 $\mathscr{M}=\widetilde{\mathscr{M}}$，则称 $\mu$ \textbf{饱和}。
\end{definition*}

\begin{enumerate}[label=\alph*.,ref=(\alph*)]
\item\label{item:ch01-sigma-finite-saturated} 若 $\mu$ 为 $\sigma$-有限测度，则 $\mu$ 饱和。
\item\label{item:ch01-locally-measurable-algebra} $\widetilde{\mathscr{M}}$ 是 $\sigma$-代数。
\item\label{item:ch01-saturation-extension} 在 $\widetilde{\mathscr{M}}$ 上定义 $\widetilde{\mu}$：若 $E\in\mathscr{M}$，则 $\widetilde{\mu}(E)=\mu(E)$；否则 $\widetilde{\mu}(E)=\infty$。则 $\widetilde{\mu}$ 是 $\widetilde{\mathscr{M}}$ 上的饱和测度，称为 $\mu$ 的\textbf{饱和化}。
\item\label{item:ch01-saturation-complete} 若 $\mu$ 完备，则 $\widetilde{\mu}$ 也完备。
\item\label{item:ch01-saturation-semifinite-extension} 设 $\mu$ 半有限。对 $E\in\widetilde{\mathscr{M}}$，定义 $\underline{\mu}(E)=\sup\{\mu(A):A\in\mathscr{M}\text{ 且 }A\subseteq E\}$。则 $\underline{\mu}$ 是 $\widetilde{\mathscr{M}}$ 上延拓 $\mu$ 的饱和测度。
\item\label{item:ch01-saturation-nonunique} 设 $X_1,X_2$ 是两两不交的不可数集，$X=X_1\cup X_2$，$\mathscr{M}$ 为 $X$ 中可数集或余可数集的 $\sigma$-代数。令 $\mu_0$ 为 $\mathscr{P}(X_1)$ 上的计数测度，并在 $\mathscr{M}$ 上定义 $\mu(E)=\mu_0(E\cap X_1)$。则 $\mu$ 为 $\mathscr{M}$ 上的测度，$\widetilde{\mathscr{M}}=\mathscr{P}(X)$；采用\subcref{exercise:ch01-saturation}{item:ch01-saturation-extension}和\subcref{exercise:ch01-saturation}{item:ch01-saturation-semifinite-extension}的记号，有 $\widetilde{\mu}\ne\underline{\mu}$。
\end{enumerate}
\end{exercise}
\end{document}
